\cleardoublepage
\newgeometry{left=1.05cm,right=1.05cm,top=1.15cm,bottom=0.95cm}

\begingroup
\setlength{\tabcolsep}{3pt}
\renewcommand{\arraystretch}{1.01}
\fontsize{10.4pt}{11.8pt}\selectfont
\sloppy
\thispagestyle{empty}

\begin{center}
    \textbf{\MakeUppercase{Đại học Bách khoa Hà Nội}}\\[-0.2ex]
    \underline{\textbf{\MakeUppercase{Trường Điện -- Điện tử}}}\\[1.5ex]
    {\fontsize{14.5pt}{16pt}\selectfont\textbf{\MakeUppercase{Phiếu đánh giá đồ án tốt nghiệp}}}\\[-0.2ex]
    \textbf{\MakeUppercase{(Dành cho cán bộ hướng dẫn)}}
\end{center}

\vspace{0.42cm}
\noindent\textbf{Tên đề tài:} \dotfill

\vspace{0.1cm}
\noindent\textbf{Sinh viên thực hiện:} \dotfill\hspace{0.7cm}\textbf{MSSV:} \dotfill

\vspace{0.1cm}
\noindent\textbf{Cán bộ hướng dẫn} \dotfill

\vspace{0.5cm}
\noindent
\begin{tabular}{|>{\centering\arraybackslash}m{1.05cm}
                |>{\centering\arraybackslash}m{2.72cm}
                |>{\setlength{\parindent}{0pt}\arraybackslash}p{12.08cm}
                |>{\centering\arraybackslash}m{1.35cm}|}
    \hline
    \textbf{STT} &
    \textbf{\shortstack{Tiêu chí\\(Điểm tối đa)}} &
    \centering\arraybackslash\textbf{Hướng dẫn đánh giá tiêu chí} &
    \textbf{\shortstack{Điểm\\tiêu\\chí}} \\
    \hline
    \multirow{2}{*}{1} &
    \multirow{2}{=}[-0.6ex]{\centering\fontsize{9.6pt}{10.5pt}\selectfont\textbf{Thái độ làm\\việc (2,5 điểm)}} &
    Nghiêm túc, tích cực và chủ động trong quá trình làm ĐATN &
    \multirow{2}{*}{} \\
    \cline{3-3}
    &
    &
    Hoàn thành đầy đủ và đúng tiến độ các nội dung được GVHD giao &
    \\
    \hline
    \multirow{2}{*}{2} &
    \multirow{2}{=}[-3.2ex]{\centering\fontsize{9.6pt}{10.5pt}\selectfont\textbf{Kỹ năng viết\\quyển (2\\điểm)}} &
    Trình bày đúng mẫu quy định, bố cục các chương logic và hợp lý: Bảng biểu, hình ảnh rõ ràng, có tiêu đề, được đánh số thứ tự và được giải thích hay đề cập đến trong đồ án, có căn lề, dấu cách sau dấu chấm, dấu phẩy, có mở đầu chương và kết luận chương, có liệt kê tài liệu tham khảo và có trích dẫn, v.v. &
    \multirow{2}{*}{} \\
    \cline{3-3}
    &
    &
    Kỹ năng diễn đạt, phân tích, giải thích, lập luận: Cấu trúc câu rõ ràng, văn phong khoa học, lập luận logic và có cơ sở, thuật ngữ chuyên ngành phù hợp, v.v. &
    \\
    \hline
    \multirow{3}{*}{3} &
    \multirow{3}{=}[-5.0ex]{\centering\fontsize{9.6pt}{10.5pt}\selectfont\textbf{Nội dung và\\kết quả (5\\điểm)}} &
    Nêu rõ tính cấp thiết, ý nghĩa khoa học và thực tiễn của đề tài, các vấn đề và các giả thuyết, phạm vi ứng dụng của đề tài. Thực hiện đầy đủ quy trình nghiên cứu: Đặt vấn đề, mục tiêu đề ra, phương pháp nghiên cứu/ giải quyết vấn đề, kết quả đạt được, đánh giá và kết luận. &
    \multirow{3}{*}{} \\
    \cline{3-3}
    &
    &
    Nội dung và kết quả được trình bày một cách logic và hợp lý, được phân tích và đánh giá thỏa đáng. Biện luận phân tích kết quả mô phỏng/ phần mềm/ thực nghiệm, so sánh kết quả đạt được với kết quả trước đó có liên quan. &
    \\
    \cline{3-3}
    &
    &
    Chỉ rõ phù hợp giữa kết quả đạt được và mục tiêu ban đầu đề ra đồng thời cung cấp lập luận để đề xuất hướng giải quyết có thể thực hiện trong tương lai. Hàm lượng khoa học/ độ phức tạp cao, có tính mới/ tính sáng tạo trong nội dung và kết quả đồ án. &
    \\
    \hline
    \multirow{2}{*}{4} &
    \multirow{2}{=}[-1.6ex]{\centering\fontsize{9.6pt}{10.5pt}\selectfont\textbf{Thành tích (1\\điểm)}} &
    Bài báo KH/ Giải SV NCKH giải 3 cấp Trường trở lên/ Giải thưởng KH/ Bằng phát minh sáng chế (1,0 điểm). &
    \multirow{2}{*}{} \\
    \cline{3-3}
    &
    &
    Báo cáo hội nghị SV NCKH/ Giải KK/ Sản phẩm ứng dụng có tính hoàn thiện cao, khối lượng thực hiện lớn (0,5 điểm). &
    \\
    \hline
    \multicolumn{3}{|r|}{\textbf{Điểm tổng các tiêu chí:}} &
    \\
    \hline
    \multicolumn{3}{|r|}{\textbf{Điểm hướng dẫn:}} &
    \\
    \hline
\end{tabular}

\vspace{0.4cm}
\begin{flushright}
    \begin{minipage}{0.35\textwidth}
        \centering
        \textbf{Cán bộ hướng dẫn}\\
        \textit{(Ký và ghi rõ họ tên)}
    \end{minipage}
\end{flushright}

\vfill
\noindent\textit{Điểm từng tiêu chí cho lẻ đến 0,5. Nếu Điểm tổng các tiêu chí \(> 10\) thì Điểm hướng dẫn làm tròn thành 10.}

\endgroup

\cleardoublepage

\begingroup
\setlength{\tabcolsep}{3pt}
\renewcommand{\arraystretch}{1.01}
\fontsize{10.4pt}{11.8pt}\selectfont
\sloppy
\thispagestyle{empty}

\begin{center}
    \textbf{\MakeUppercase{Đại học Bách khoa Hà Nội}}\\[-0.2ex]
    \underline{\textbf{\MakeUppercase{Trường Điện -- Điện tử}}}\\[1.5ex]
    {\fontsize{14.5pt}{16pt}\selectfont\textbf{\MakeUppercase{Phiếu đánh giá đồ án tốt nghiệp}}}\\[-0.2ex]
    \textbf{\MakeUppercase{(Dành cho cán bộ phản biện)}}
\end{center}

\vspace{0.42cm}
\noindent\textbf{Tên đề tài:} \dotfill

\vspace{0.1cm}
\noindent\textbf{Sinh viên thực hiện:} \dotfill\hspace{0.7cm}\textbf{MSSV:} \dotfill

\vspace{0.1cm}
\noindent\textbf{Cán bộ hướng dẫn} \dotfill

\vspace{0.5cm}
\noindent
\begin{tabular}{|>{\centering\arraybackslash}m{1.05cm}
                |>{\centering\arraybackslash}m{2.72cm}
                |>{\setlength{\parindent}{0pt}\arraybackslash}p{12.08cm}
                |>{\centering\arraybackslash}m{1.35cm}|}
    \hline
    \textbf{STT} &
    \textbf{\shortstack{Tiêu chí\\(Điểm tối đa)}} &
    \centering\arraybackslash\textbf{Hướng dẫn đánh giá tiêu chí} &
    \textbf{\shortstack{Điểm\\tiêu\\chí}} \\
    \hline
    \multirow{2}{*}{1} &
    \multirow{2}{=}[-3.2ex]{\centering\fontsize{9.6pt}{10.5pt}\selectfont\textbf{Trình bày\\quyển ĐATN\\(4 điểm)}} &
    Đồ án trình bày đúng mẫu quy định, bố cục các chương logic và hợp lý: Bảng biểu, hình ảnh rõ ràng, có tiêu đề, được đánh số thứ tự và được giải thích hay đề cập đến trong đồ án, có căn lề, dấu cách sau dấu chấm, dấu phẩy, có mở đầu chương và kết luận chương, có liệt kê tài liệu tham khảo và có trích dẫn, v.v. &
    \multirow{2}{*}{} \\
    \cline{3-3}
    &
    &
    Kỹ năng diễn đạt, phân tích, giải thích, lập luận: cấu trúc câu rõ ràng, văn phong khoa học, lập luận logic và có cơ sở, thuật ngữ chuyên ngành phù hợp, v.v. &
    \\
    \hline
    \multirow{3}{*}{2} &
    \multirow{3}{=}[-5.0ex]{\centering\fontsize{9.6pt}{10.5pt}\selectfont\textbf{Nội dung và\\kết quả đạt\\được (5,5)\\điểm}} &
    Nêu rõ tính cấp thiết, ý nghĩa khoa học và thực tiễn của đề tài, các vấn đề và các giả thuyết, phạm vi ứng dụng của đề tài. Thực hiện đầy đủ quy trình nghiên cứu: Đặt vấn đề, mục tiêu đề ra, phương pháp nghiên cứu/ giải quyết vấn đề, kết quả đạt được, đánh giá và kết luận. &
    \multirow{3}{*}{} \\
    \cline{3-3}
    &
    &
    Nội dung và kết quả được trình bày một cách logic và hợp lý, được phân tích và đánh giá thỏa đáng. Biện luận phân tích kết quả mô phỏng/ phần mềm/ thực nghiệm, so sánh kết quả đạt được với kết quả trước đó có liên quan. &
    \\
    \cline{3-3}
    &
    &
    Chỉ rõ phù hợp giữa kết quả đạt được và mục tiêu ban đầu đề ra đồng thời cung cấp lập luận để đề xuất hướng giải quyết có thể thực hiện trong tương lai. Hàm lượng khoa học/ độ phức tạp cao, có tính mới/ tính sáng tạo trong nội dung và kết quả đồ án. &
    \\
    \hline
    \multirow{2}{*}{3} &
    \multirow{2}{=}[-2.0ex]{\centering\fontsize{9.6pt}{10.5pt}\selectfont\textbf{Điểm thành\\tích (1 điểm)}} &
    Có bài báo KH được đăng hoặc chấp nhận đăng/ đạt giải SV NCKH giải 3 cấp Trường trở lên/ các giải thưởng KH quốc tế/ trong nước từ giải 3 trở lên/ Có đăng ký bằng phát minh sáng chế. (1 điểm) &
    \multirow{2}{*}{} \\
    \cline{3-3}
    &
    &
    Được báo cáo tại hội đồng cấp Trường trong hội nghị SV NCKH nhưng không đạt giải từ giải 3 trở lên/ Đạt giải khuyến khích trong các kỳ thi quốc gia và quốc tế khác về chuyên ngành. (0,5 điểm) &
    \\
    \hline
    \multicolumn{3}{|r|}{\textbf{Điểm tổng các tiêu chí:}} &
    \\
    \hline
    \multicolumn{3}{|r|}{\textbf{Điểm phản biện:}} &
    \\
    \hline
\end{tabular}

\vspace{0.4cm}
\begin{flushright}
    \begin{minipage}{0.35\textwidth}
        \centering
        \textbf{Cán bộ phản biện}\\
        \textit{(Ký và ghi rõ họ tên)}
    \end{minipage}
\end{flushright}

\vfill
\noindent\textit{Điểm từng tiêu chí cho lẻ đến 0,5. Điểm phản biện tối đa là 10.}

\endgroup
\cleardoublepage
\restoregeometry
